% Generated draft: requires mathematical and source review.
\documentclass[11pt]{article}
\usepackage[margin=25mm]{geometry}
\usepackage{fontspec}
\setmainfont{Latin Modern Roman}
\usepackage{amsmath,amssymb,mathtools,unicode-math}
\setmathfont{Latin Modern Math}
\usepackage{enumitem,needspace,xurl,hyperref}
\setlength{\parindent}{0pt}
\setlength{\parskip}{5pt}
\setlength{\emergencystretch}{3em}
\newcommand{\N}{\mathbb{N}}
\newcommand{\Z}{\mathbb{Z}}
\newcommand{\Q}{\mathbb{Q}}
\newcommand{\R}{\mathbb{R}}
\newcommand{\C}{\mathbb{C}}
\newcommand{\F}{\mathbb{F}}
\newcommand{\A}{\mathbb{A}}
\newcommand{\PP}{\mathbb{P}}
\newcommand{\E}{\mathbb{E}}
\newcommand{\eps}{\varepsilon}
\newcommand{\dd}{\,\mathrm{d}}
\begin{document}
\section*{UnsolvedMath ID: 2307044}
\textbf{Source problem code:} catalogue entry
\subsection{Definitions and mathematical statement}
For \(\beta>0\), let \(\mathcal H^\beta\) be \(\beta\)-dimensional Hausdorff measure on \(\R^2\cong\C\). A finite Borel measure \(\nu\) is singular with respect to \(\mathcal H^\beta\) when there is a Borel set \(E\) with
\[
\nu(\R^2\setminus E)=0,\qquad \mathcal H^\beta(E)=0.
\]
Let \(U\subseteq\C\) be open and \(a\in U\); write \(\lambda_a^U\) for harmonic measure at \(a\), using the component of \(U\) containing \(a\). In its usual domain setting, harmonic measure represents boundary values of harmonic functions, or equivalently the Brownian exit distribution when that interpretation applies. The first question is
\[
\forall U\ \forall a\in U\ \forall\beta>1,\qquad
\lambda_a^U\perp\mathcal H^\beta.
\]
The source credits \O{}ksendal with \(\lambda_a^U\) singular relative to planar area measure and asks for the stronger Hausdorff-measure statement. The nontrivial new exponent range is \(1<\beta<2\); the question retains all \(\beta>1\).

The second question asks the same singularity for the Keldysh measure associated with a compact set \(K\subseteq\C\) at a point \(a\in K\):
\[
\forall K\text{ compact}\ \forall a\in K\ \forall\beta>1,\qquad
\mu_a^K\perp\mathcal H^\beta.
\]
Here \(\mu_a^K\) denotes the source's intended Keldysh measure, whose construction and representing class are not supplied. The source actually writes \(\mu_a^U\) while saying ``with respect to a compact set \(K\)''; the superscript \(K\) in this display is an explicitly proposed notational identification. The source also credits \O{}ksendal with its area-measure singularity. This second measure is not silently identified with ordinary harmonic measure on an open domain or on the complement of \(K\).

The supplied OPEN-TRIAGE summary reports open status in the 2018 edition and no definitive resolution located, without distinguishing a precise current status for the two measures. This is supplied metadata rather than a verified literature assessment.
\subsection{Short English statement}
Are planar harmonic measure and the compact-set Keldysh measure carried by Hausdorff-null sets for every dimension exponent above one? The Keldysh definition and notation, and the current status of each part, need source review.
\subsection{Sources}
\begin{itemize}
\item Local UnsolvedMath registry, permanent ID \(2307044\), exact source code \texttt{catalogue entry}, ``Research Problems in Function Theory --- Problem 7.44.''
\item Hayman, \emph{Research Problems in Function Theory} (2018), Problem 7.44, attributed to B. \O{}ksendal and crediting his area-singularity results; \url{https://arxiv.org/abs/1809.07200}. Source-TeX extraction; recorded access 2026-07-29.
\item The \O{}ksendal results and Keldysh construction have no complete references in the supplied excerpt. The OPEN-TRIAGE summary supplies no separate authoritative current-status citation.
\end{itemize}
\end{document}
